% Multiplication by e^{i theta} rotates a complex number through the angle
% theta (counter-clockwise for theta > 0) without changing its length.
\documentclass{article}
\usepackage[paperwidth=120cm,paperheight=120cm,margin=10mm]{geometry}
\pagestyle{empty}
\usepackage{fontspec}
\usepackage[bidi=basic]{babel}
\babelprovide[import]{english}
\babelprovide[main, import]{persian}
\babelfont{rm}[Renderer=Node, Path=../fonts/, Extension=.ttf, UprightFont=*-Regular, BoldFont=*-Bold]{Vazirmatn}
\babelfont{sf}[Renderer=Node, Path=../fonts/, Extension=.ttf, UprightFont=*-Regular, BoldFont=*-Bold]{Vazirmatn}
\providecommand{\setRTL}{}\providecommand{\setLTR}{}
\providecommand{\RL}[1]{\foreignlanguage{persian}{#1}}
\providecommand{\LR}[1]{\babelsublr{#1}}
\usepackage{tikz}
\usepackage{pgfplots}
\pgfplotsset{compat=1.18}
\usetikzlibrary{arrows.meta, positioning, shapes.geometric, shapes.misc, calc, fit, backgrounds, decorations.pathmorphing, patterns, angles, quotes}
\begin{document}
\setRTL
\babelsublr{\begin{tikzpicture}[>=Stealth, font=\small, scale=1.5]
  % axes
  \draw[->, thick] (-2.2,0) -- (2.4,0) node[right] {$\mathrm{Re}$};
  \draw[->, thick] (0,-0.5) -- (0,2.4) node[above] {$\mathrm{Im}$};

  \coordinate (O) at (0,0);
  % z1 at 20 degrees, radius 2
  \coordinate (Z1) at (20:2);
  % z2 = e^{i theta} z1, rotated by theta=60 -> at 80 degrees, same radius
  \coordinate (Z2) at (80:2);

  % unit-radius reference arc to show equal length is preserved
  \draw[dashed, gray] (20:2) arc (20:80:2);

  % original vector z1
  \draw[->, very thick, blue!70!black] (O) -- (Z1) node[right] {$z_1 = r\,e^{i\theta_1}$};
  % rotated vector z2
  \draw[->, very thick, red!75!black] (O) -- (Z2) node[above right] {$z_2 = e^{i\theta}\,z_1$};

  % rotation angle theta between them
  \draw[->, thick] (50:0.95) arc (20:80:0.95);
  \node at (50:1.3) {$\theta$};

  % note equal magnitudes
  \node[align=center, gray] at (1.7,-0.42) {طول یکسان $r$\\(فقط چرخش)};
\end{tikzpicture}}
\end{document}
